\section{Related Work}
\label{sec:related}

\paragraph{Video retrieval augmented generation.}
RAG grounds generation in retrieved evidence. For long videos, Video-RAG~\cite{luo2024videorag} augments a VLM with
visually aligned auxiliary texts, such as speech transcripts, on-screen text, and detected objects, retrieved for each
query, and VideoRAG~\cite{jeong2025videorag} retrieves whole videos from a corpus as external knowledge. A second line
of work structures the index. VideoRAG~\cite{ren2025videorag} builds a graph-based knowledge index across videos,
Vgent~\cite{shen2025vgent} represents a video as a graph of clips that preserves their semantic and temporal
dependencies, followed by a structured reasoning step that filters retrieval noise, and AVA~\cite{yan2025ava} indexes a
video as an event knowledge graph that an agent searches. Other systems add spatio-temporal
structure~\cite{fu2026videostir} or temporal alignment~\cite{cao2025tvrag} to retrieval, and
\citet{lee2026rethinking} benchmark retrieval and generation separately.
These systems, however, focus on what to index and retrieve, and pass every retrieved item to the VLM at its full
cost. \name instead decides how many tokens each retrieved item receives, and can be placed after any of these
retrievers.

\paragraph{Keyframe selection.}
Most VLMs sample frames uniformly and therefore often miss the content that a query needs. Query-aware methods instead
select frames by jointly optimizing their relevance and temporal coverage~\cite{tang2025aks,peng2026qca}, by greedy
marginal utility~\cite{shi2026markey}, by learning to query informative combinations of
frames~\cite{yu2024framevoyager}, with the attention of a small VLM used as a selector~\cite{wang2026dafs}, or by
exploring clips as the arms of a bandit~\cite{zhu2026focus}. All of these methods, however, return a subset of frames,
each at the same token cost. \name instead treats the token count of each frame as a decision variable.

\paragraph{Token allocation across frames.}
Closest to our work are methods that vary the token count of frames. Q-Frame~\cite{zhang2025qframe} shows frames at
three resolutions according to their relevance rank, QuoTA~\cite{luo2025quota} assigns tokens in proportion to a
query-oriented frame score, and AdaptToken~\cite{qi2026adapttoken} distributes a global token budget over groups of
frames according to the response entropy of the VLM. Concurrent with our work, AdaAlloc~\cite{an2026adaalloc} lets the
VLM plan how to split the budget between sparse low-resolution frames and dense high-resolution frames; it became
public after the CVPR~2027 cutoff for prior work, and we include it in our comparisons for completeness. All of these
methods allocate with a fixed rule or a heuristic plan. \name instead derives the allocation as the solution of an
optimization problem with a learned token value, extends it to transcripts, and maintains it across the turns of a
conversation.

\paragraph{Token compression and KV-cache reuse.}
A complementary line of work reduces the tokens of a given input. FastV~\cite{chen2024fastv} and
VisionZip~\cite{yang2024visionzip} prune visual tokens by attention, ToMe~\cite{bolya2023tome} merges similar tokens
inside the vision encoder, and LongVU~\cite{shen2024longvu}, DyCoke~\cite{tao2024dycoke}, and
MARC~\cite{wu2026marc} compress video tokens by exploiting their temporal redundancy. For language models,
H$_2$O~\cite{zhang2023h2o} and StreamingLLM~\cite{xiao2024streamingllm} evict entries of the KV cache, and
ReKV~\cite{rekv2025} retrieves cached video KV entries for streaming question answering. These methods, however,
compress each frame or cache independently. \name uses a pruner only to realize a token count, decides the counts
jointly, and reuses cached tokens instead of evicting them.
